\section{Scope, conventions, and the shape of the main theorem}
\label{sec:scope}

\subsection{What is proved}

Throughout, $k\subseteq K$ are algebraically closed fields and $K$ has infinite
transcendence degree over $k$ (\cref{conv:ud}).  All tuples are finite tuples of
elements of $K$; $\acl$ is relative algebraic closure in $K$ over $k$, and the
rank $\rk(a/B)$ of a tuple $a$ over a set $B$ is the transcendence degree of
$k(B,a)$ over $k(B)$ (\cref{sec:ud}).

The object of study is a \emph{group configuration} (\cref{def:gc}): six tuples
$a,b,c,x,y,z$ with the four \emph{lines}
\[
  \{a,b,c\},\qquad \{a,x,y\},\qquad \{b,y,z\},\qquad \{c,z,x\},
\]
such that (i) each group point is algebraic over the two others on the group line,
(ii) on each space line the two space points are interalgebraic over the group point,
(iii) every non-collinear triple is independent, and (iv) each group point is
interalgebraic with the \emph{canonical base} of the correspondence it defines
between the two space points on its line (the ``no smaller field of definition''
condition; in the reconstruction blueprint it is clause (iv) of $\Psi$).  Group
points have a common rank $n\ge 1$, space points a common rank $m$, $1\le m\le n$.

\begin{theorem}[main theorem, informal; precise form in \cref{thm:main}]
\label{thm:main-informal}
Let $(a,b,c,x,y,z)$ be a group configuration over $k$.  Then there are finitely many
elements $e$ of $K$, independent from $abcxyz$ over $k$, such that, writing
$k'=\acl(k,e)$ for the new base field, there exist
\begin{itemize}[nosep]
\item a group $G$ of germs of quasi-rational self-maps of an irreducible $k'$-variety
  $V$ of dimension $m$, connected, of rank $n$, whose elements of rank $n$ are exactly
  the germs $f_1^{-1}f_2$ for independent members $f_1,f_2$ of a generic set $\mathcal F$
  of invertible germs from $V$ to a second variety $V'$ (a single $\Gamma$-orbit, hence
  canonically parametrised by the generic points of one variety of codes);
\item a transitive faithful $G$-set $S$ containing the generic points of $V$;
\item generic independent $g,h\in G$ and a generic $s\in V(K)$,
\end{itemize}
such that, with $x'=g\cdot s$, $y'=s$, $z'=h\cdot s$, $a'=g$, $b'=h$, $c'=hg^{-1}$,
each of $a,b,c,x,y,z$ is interalgebraic over $k'$ with the corresponding primed
element (group elements being represented by their canonical codes, which are tuples
of $K$).  Moreover, for a generic point $y_0$ of $V$ over $k'(abcxyz)$ and a further
tuple $t$ of fresh generic points (chosen after $y_0$) with $(e,t)$ independent from
$abcxyz$, $G$ is \emph{constructible over $k'(t)$} (\cref{def:constructible-group}) and
$S$ is a constructible $G$-set over $k'(t,y_0)$, namely $G/H$ for the constructible
stabiliser $H$ of $\iota(y_0)$ (\cref{thm:constructible}).
\end{theorem}

In the application to the reconstruction blueprint (\cref{sec:application}),
$n=2$, $m=1$, and the configuration is the $\Psi$-witness of the geometric
relation $Q$; the theorem supplies the two-dimensional group acting on a curve whose
classification (blueprint Lemma \texttt{affine-action}) yields the affine coordinates
of Lemma \texttt{affine-grid-extraction}.  The classification itself is not part of
this document.

\subsection{What is not proved here}

\begin{itemize}
\item \textbf{Weil's theorem.}  We do not prove that a constructible group over an
  algebraically closed field is (isomorphic to) an algebraic group in the sense of
  algebraic geometry, nor do we construct a variety structure on $G$.  The main
  theorem is an ``up to interalgebraicity'' statement and every quantity in it is
  defined in terms of fields, ideals of polynomials and constructible sets.
  \Cref{sec:constructibility} records the precise definable-group structure that
  is obtained, which is the input any later algebraization would start from.
\item \textbf{Classification} of the resulting group (dimension two acting on a
  curve, blueprint Lemma \texttt{affine-action}).
\item \textbf{Elimination of imaginaries in full generality.}  Only the codes that
  are actually needed are constructed explicitly (\cref{sec:codes}); the
  constructible group of \cref{sec:constructibility} is presented as a quotient of a
  constructible set by a constructible equivalence relation, and we do not prove
  that such a quotient can be identified with a constructible subset of some $K^N$.
\end{itemize}

\subsection{How the proof is organised}

\begin{enumerate}
\item \Cref{sec:ud}: the universal domain, rank calculus, automorphisms, loci.
\item \Cref{sec:fod}: smallest fields of definition of ideals and the resulting
  \emph{canonical bases} $\Cb(a/L)$; their Galois and independence properties.
\item \Cref{sec:codes}: definable closure as perfect closure; explicit codes of
  finite sets of tuples.
\item \Cref{sec:germs}: germs of quasi-rational maps between irreducible varieties,
  their evaluation at generic points, composition, inverses, and the dictionary
  between a triple $(b,x,y)$ with $y\in\dcl(b,x)$ and a germ whose code is
  $\Cb(xy/b)$.
\item \Cref{sec:hw}: the Hrushovski--Weil theorem: a generically transitive generic
  set of invertible germs (a single orbit under the automorphism group) closed under
  inverse and generic composition generates a connected group of germs whose elements
  of maximal rank are exactly that set, and a transitive faithful homogeneous space.
\item \Cref{sec:bij}: the version for families of bijections between two varieties
  (composites $f_{b_1}^{-1}f_{b_2}$), with the rank criterion for the independence
  hypothesis.
\item \Cref{sec:gc}: group configurations, their invariance under base change and
  interalgebraic replacement, and the converse (configurations of a homogeneous space).
\item \Cref{sec:proof}: the three steps of Hrushovski's proof (making the lines
  definable, the independence hypothesis, identification) and the main theorem.
\item \Cref{sec:constructibility}: the group and the space are constructible
  (definable) over the extended base.
\item \Cref{sec:descent}: descent of the extra parameters into $k$ for first-order
  consequences.
\item \Cref{sec:application}: the $\Psi$-witness is a group configuration; what the
  theorem gives for the reconstruction blueprint.
\item \Cref{sec:lean}: proposed Lean architecture and the Mathlib/AclGeom API used.
\end{enumerate}

\subsection{Conventions}

\begin{convention}[orientation]\label{conv:orientation}
We use the orientation of Bays' notes, which coincides with the letters of the
reconstruction blueprint's $\Psi$-witness: the group line is $\{a,b,c\}$ with
\emph{$b=c\ast a$}, and the space lines read
\[
   x=a\ast y,\qquad z=b\ast y,\qquad z=c\ast x .
\]
Thus $a$ carries $y$ to $x$, $c$ carries $x$ to $z$, and $b=c\ast a$ carries $y$ to
$z$.  The independent triple of the configuration is $(a,b,y)$.  In the
$\Psi$-witness of the reconstruction blueprint, $(a,b,c,x,y,z)=(A,B,C,X,Y,Z)$:
$X\le A\vee Y$, $Z\le B\vee Y$, $Z\le C\vee X$ (\cref{sec:application}).
\end{convention}

\begin{convention}[citations to Lean]\label{conv:lean}
Statements that are already available in Mathlib (revision
\lean{d13f23b7}, tag \lean{v4.34.1}) or in the project library \lean{AclGeom} are
marked with their declaration names in \lean{teletype}.  Such statements are
\emph{not} reproved here.  Every other lemma is proved in full.  The Lean names
used are listed in \cref{sec:lean}; when a cited name has a slightly different
form from the statement in the text (for instance a lattice formulation versus a
tuple formulation), the required bridge is stated explicitly there.
\end{convention}

\begin{convention}[characteristic]\label{conv:char}
The characteristic $p\ge 0$ of $k$ is arbitrary.  In characteristic $p>0$, the
definable closure of a field is its perfect closure and ``definable functions'' are
quasi-rational (rational after extracting $p$-power roots).  No separability
hypothesis is ever used.  All constructions are uniform in $p$; in characteristic
$0$ every ``quasi-rational'' map is rational.
\end{convention}
